\par\noindent\begin{minipage}{\linewidth}
% !TeX root = ../教材题库_全部条目_转写工作稿.tex
% Source: XJ-B1-C01-PR-P019-Y0389
% 第1章 / 1.2 常用逻辑用语 / 教材 p.19 / PDF 第26页
% 已对照原页转写题干及公式；待教师复核；未制作答案
% 原题及相邻原解截图（核对用）：../题目原图/XJ-B1-C01-PR-P019-Y0389-01.webp
% 本文件只在转写完成后作为题干使用；OCR 条目尚不是可靠题干。
\begin{bankitem}{p.19 练习 第2题}{第1章 / 1.2 常用逻辑用语 / 教材 p.19 / PDF 第26页}{XJ-B1-C01-PR-P019-Y0389}{已对照原页转写题干及公式；待教师复核；未制作答案}
从“充分而不必要条件”“必要而不充分条件”“充要条件”与“既不充分又不必要条件”中选出适当的一种填空：
\subq{1}$\triangle ABC$ 中，$\angle C=90^\circ$ 是 $AB^2=AC^2+BC^2$ 的\blank；
\subq{2}$x>0$ 是 $x\ge1$ 的\blank；
\subq{3}$x=2$ 是 $x^2=4$ 的\blank；
\subq{4}$0<x<2$ 是 $1<x<3$ 的\blank。
\end{bankitem}
\end{minipage}\par\medskip
